\section{Auswirkungen eines Fehlerfalls auf Schutzziele}

Die Auswirkungen der beiden untersuchten Szenarien sind in Tabelle~\ref{tab:Auswirkungen-Schutzziele} im Anhang (Abschnitt~\ref{app:auswirkungstabelle}) zusammengefasst.

\subsection*{Angriff POODLE} 
Der POODLE-Angriff (Padding Oracle On Downgraded Legacy Encryption, CVE-2014-3566) zeigt, wie ein Angreifer eine Verbindung von TLS auf SSL~3.0 downgraded und dadurch die Verbindung kompromittiert \cite{cve-poodle}. Die Padding-Oracle-Eigenschaft von SSL~3.0 im CBC-Modus ermöglicht die iterative Rekonstruktion von Klartextbytes durch einen Angreifer im Netzpfad (\textit{Vertraulichkeit}). Manipulierte CBC-Blöcke werden akzeptiert, sofern das Padding zufällig korrekt ist -- die MAC-then-Encrypt-Konstruktion bietet keinen ausreichenden Schutz (\textit{Integrität}). Mit dem erzwungenen Wechsel auf SSL~3.0 entfallen die Finished-Message-Signaturen moderner TLS-Versionen; ohne \texttt{TLS\_FALLBACK\_SCSV} bleibt der Downgrade unbemerkt (\textit{Authentizität}) \cite{thomas-krenn-poodle}.

\subsection*{Selbstsigniertes Zertifikat} 

Präsentiert ein Server ein selbstsigniertes Zertifikat, bricht ein korrekt konfigurierter Client den Handshake mit einem \texttt{CERTIFICATE\_UNKNOWN}-Alert ab -- die Verbindung kommt nicht zustande \cite{mdndocs}. Deaktiviert ein Nutzer diese Prüfung jedoch (z.\,B.\ durch Akzeptieren einer Browser-Warnung oder mittels \texttt{-{}-no-verify}-Flag), öffnet sich ein Angriffsfenster: Ein Angreifer im Netzpfad kann ein eigenes selbstsigniertes Zertifikat einschleusen, das der Client mangels CA-Prüfung ebenfalls akzeptiert. Damit terminiert der Angreifer \emph{beide} TLS-Verbindungen separat -- zum Client und zum echten Server -- und leitet Daten transparent weiter (\textit{Man-in-the-Middle}). Alle übertragenen Klartextdaten sind dem Angreifer zugänglich, obwohl der Nutzer eine \texttt{https://}-Verbindung im Browser sieht (\textit{Vertraulichkeit verletzt, Authentizität verletzt}). Die TLS-Record-MAC-Prüfung schützt die Integrität \emph{innerhalb} jeder der beiden Teilverbindungen -- der Angreifer kann jedoch Inhalte modifizieren, bevor er sie in der zweiten Verbindung weitersendet, sodass der Endnutzer manipulierte Daten als integer wahrnimmt (\textit{Integrität de facto verletzt}).

\section{Härtungsmaßnahmen}
\subsection*{TLS Versionsbeschränkung}

Das BSI schreibt in der \textbf{Technischen Richtlinie TR-02102-2} (Stand: Januar 2026) vor, dass in neuen Systemen ausschließlich TLS~1.2 und TLS~1.3 eingesetzt werden sollen \cite{bsi-tr02102}. SSL~2.0, SSL~3.0 und TLS~1.0 sind vollständig zu deaktivieren; TLS~1.1 wird ebenfalls nicht mehr empfohlen. Die Konfiguration in nginx lautet entsprechend:

\begin{lstlisting}
ssl_protocols TLSv1.2 TLSv1.3;
\end{lstlisting}

Durch die Abschaltung aller Legacy-Protokolle entfällt die Angriffsfläche für POODLE und BEAST \cite{cve-beast} grundlegend, da kein Protokoll-Fallback mehr stattfinden kann \cite{thomas-krenn-poodle}. Ergänzend sollte \textbf{TLS\_FALLBACK\_SCSV} (RFC~7507) implementiert werden: Übermittelt ein Client diesen Signaling Cipher Suite Value in einem erneuten, herabgestuften Verbindungsversuch, bricht der Server den Handshake mit einem \texttt{inappropriate\_fallback}-Alert ab und verhindert so auch aktive Downgrade-Versuche bei Systemen, die Legacy-Versionen aus Kompatibilitätsgründen noch führen \cite{rfc7507}.

\subsection*{Sichere Cipher Suites}

Algorithmen ohne kryptographische Langzeitsicherheit -- darunter RC4, DES, 3DES, Export-Cipher und MD5-basierte MACs -- sind vollständig aus der Konfiguration zu entfernen \cite{bsi-tr02102}. Gemäß BSI~TR-02102-2 sind ausschließlich Cipher Suites mit \textbf{Perfect Forward Secrecy (PFS)} zu verwenden: Nur ephemere Schlüsselaustauschverfahren (ECDHE, DHE) stellen sicher, dass ein nachträglich kompromittierter Langzeitschlüssel keine bereits aufgezeichneten Sessions entschlüsselbar macht. Empfohlene Suiten für TLS~1.2 sind:

\begin{lstlisting}
TLS_ECDHE_RSA_WITH_AES_256_GCM_SHA384
TLS_ECDHE_ECDSA_WITH_AES_128_GCM_SHA256
TLS_DHE_RSA_WITH_AES_256_GCM_SHA384
\end{lstlisting}

Für TLS~1.3 sind die verfügbaren Cipher Suites protokollseitig bereits auf AEAD-Verfahren (AES-GCM, ChaCha20-Poly1305) beschränkt; RC4, CBC sowie der statische RSA-Schlüsseltransport sind im Standard nicht mehr vorgesehen \cite{windows-papst-ciphers,mozilla-ssl-gen}.

\subsection*{Zertifikatsvalidierung und Vertrauensketten}

Selbstsignierte Zertifikate sind in produktiven Umgebungen unzulässig. Es sind Zertifikate einer anerkannten Zertifizierungsstelle einzusetzen, deren Root-Zertifikate in den Betriebssystem- und Browser-Trust-Stores enthalten sind. Zur weiteren Absicherung der Vertrauenskette empfehlen sich folgende Mechanismen:

\begin{itemize}
  \item \textbf{Certificate Transparency (CT):} 
    Alle ausgestellten Zertifikate werden in öffentliche, append-only-Logs eingetragen. Clients und Monitore können überprüfen, ob ein präsentiertes Zertifikat dort verzeichnet ist, und unautorisierte Ausstellungen frühzeitig erkennen.
  \item \textbf{OCSP Stapling:} 
    Der Server liefert den aktuellen Sperrstatus seines Zertifikats unmittelbar im TLS-Handshake mit, sodass der Client keinen separaten OCSP-Request senden muss. Dies     verhindert, dass ein Angreifer OCSP-Abfragen blockiert und ein bereits widerrufenes Zertifikat weiterhin nutzbar bleibt.
  \item \textbf{CAA-DNS-Einträge:} 
    \emph{Certification Authority Authorization} Records schränken auf DNS-Ebene ein, welche Zertifizierungsstellen für eine Domain überhaupt Zertifikate ausstellen dürfen, und verringern so das Risiko unberechtigter Ausstellungen.
\end{itemize}

\subsection*{Regelmäßige Überprüfung und Patch-Management}

Da Schwachstellen kontinuierlich neu entdeckt werden (vgl. MITRE CVE-Datenbank), ist eine \textbf{regelmäßige Überprüfung der TLS-Konfiguration} zwingend erforderlich. Werkzeuge wie \href{https://testssl.sh}{testssl.sh} oder der \href{https://www.ssllabs.com/ssltest/}{Qualys SSL~Labs Server Test} erlauben eine automatisierte Analyse aktiver Protokollversionen, Cipher Suites und Zertifikatseigenschaften. Das BSI macht die TR-02102-2 für Bundesbehörden gemäß § 8 Abs.~1 BSIG verbindlich und aktualisiert die Richtlinie fortlaufend -- zuletzt im Januar 2026 \cite{kes-bsi-tr}.
